\cotefile{157-3}
\scanpage{2}
\uline{Retour sur les $\unsure{MW|M\mathcal W}^{-1}$}

(\unsure{comparaison|Comparaison} avec Quillen \unsure{etc})

(IX)
% printed form
% Reverse-side print visible through the cover, read in its own direction. Right edge cropped.
{\ttfamily IEBCOPY MESSAGES AND \unsure{CONTRO|CONTROL}\\
COPY OUTDD=VERS,INDD=DE\\
\unsure{IEB139I} \illegible{3} LOAD \unsure{-} \illegible{1} \unsure{PREDONTA},\illegible{1},486,TA,\unsure{D}\\
IEB151I JOB HAS TERMINATED WITH ERROR(S)\\
\unsure{IEB147I} END OF JOB\unsure{---08|---0} WAS HIGHEST SEVERITY CODE}
% archive sheet number: 1

\scanpage{3}\mspage{1}
\deleted{\unsure{Les}} Retour \deleted{\unsure{sur|aux}}\unsure{sur} les catégories de fractions $M\Sigma^{-1}$.

1. \uline{Fonctorialité du} $\mathrm{Ch}^{\Sigma}_{\infty}(a,b)$ \uline{en} $a,b$ \uline{dans} $M\Sigma^{-1}$

On se donne une cat. $M$ avec ens. admissible de flèches $\Sigma\subset\mathrm{Fl}(M)$, d'où (en prenant p.ex. $\underline t_{\infty}=\bigcup\underline t_n$), avec $\underline t_n=\deleted{\unsure{\xi}}(\underbrace{\Delta_+\Delta_-}_{\tau_0})^{2n}$, et $\underset{\substack{\Vert\\\tau_0^{2n+2}}}{t_{n+1}}\unsure{\longrightarrow|\text{---}}\underset{\substack{\Vert\\\tau_0^{2n}}}{t_n}$ défini \deleted{par les} \deleted{\unsure{composé}} $p\unsure{*|\times}\mathrm{id}_{\tau_0^{2n}}\unsure{\times|*}p$, où $p:\tau_0\to1$ la dégénérescence. On a alors une théorie des $\mathrm{Ch}_{\Sigma;\infty}(M,a;b)$, que je note simplement $\mathrm{Ch}_{\deleted{\unsure{M}}\infty}(a,b)$\inserted{\unsure{en} \unsure{tant|vus} \unsure{qu'}objets de $\mathrm{Hot}$}. Des \deleted{\unsure{lois}} accouplements

\[
\mathrm{Ch}_{\deleted{\unsure{M}}\infty}(a,b)\times\mathrm{Ch}_{\deleted{\unsure{M}}\infty}(b,c)\longrightarrow\mathrm{Ch}_{\deleted{\unsure{M}}\infty}(a,c)\qquad\text{dans }\mathrm{Hot}
\]
avec leurs \unsure{liens} ($=$ admettant) \uline{associativité} \deleted{\unsure{en}} dans $\mathrm{Hot}$. \deleted{\unsure{Enfin, on a}} Cela définit donc des accouplements sur les $\pi_0(\mathrm{Ch}_{\deleted{\unsure{M}}\infty}(a,b))$\inserted{$\overset{\text{déf}}{=}\pi_0(a,b)$}, qui \unsure{n'est|n'est rien} autre que la composition des \deleted{flèches} \deleted{dans} $M\Sigma^{-1}$. D'autre part on a pour $a,b\in\mathrm{Ob}M$

\[
\mathrm{Hom}_M(a,b)\longrightarrow\pi_0(\mathrm{Ch}_{\infty}(a,b))=\pi_0(a,b)
\]
\unsure{compatible|compatibles} avec les compositions. Mais on peut \unsure{l'ignorer} pour l'instant, et regarder $\mathrm{Ob}M$, \inserted{ou simplement $M_0$,} avec la seule structure
\[
\left\{\begin{aligned}
M_0\times M_0&\longrightarrow\mathrm{Hot}&(a,b)&\unsure{\mapsto|\vdash}\mathrm{Ch}_{\infty}(a,b)\\
M_0\times M_0\times M_0&\longrightarrow\mathrm{Fl}\,\mathrm{Hot}&(a,\underline b,c)&\mapsto\bigl(\mathrm{Ch}_{\infty}(a,b)\times\mathrm{Ch}_{\infty}(b,c)\\
&&&\hspace{2em}\underset{\text{fl. dans Hot}}{\unsure{\twoheadrightarrow|\longrightarrow}}\mathrm{Ch}_{\infty}(a,c)\bigr)
\end{aligned}\right.
\]
avec \inserted{\unsure{condition}} associativité, et aussi une condition d'unités : $\mathrm{Ch}_{\infty}(a,a)$ est \unsure{muni|munie} d'un
\[
1_a\in\pi_0\,\mathrm{Ch}_{\infty}(a,a)=\pi_0(a,a),\quad\deleted{\unsure{\text{pour}}}
\]
\[
\deleted{\unsure{\pi_0(a,b)=\pi_0\mathrm{Ch}_0(a,b)}}
\]
jouant le rôle d'une ``identité bilatère'', dans le sens suivant. Notons que les accouplements
\[
\begin{aligned}
\mathrm{Ch}_{\infty}(a,b)\times\mathrm{Ch}_{\infty}(b,c)&\longrightarrow\mathrm{Ch}_{\infty}(a,c)\\
P\times Q&\longrightarrow R
\end{aligned}
\]
% archive sheet number: 2

\scanpage{4}\mspage{2}
définissent
\[
\begin{aligned}
P&\longrightarrow\underline{\mathrm{Hom}}_{\mathrm{Hot}}(Q,R)\\
Q&\unsure{\text{---}|\longrightarrow}\underline{\mathrm{Hom}}_{\mathrm{Hot}}(P,R)
\end{aligned}
\]
d'où en passant aux $\pi_0$
\[
\begin{aligned}
\pi_0(P)&\longrightarrow\mathrm{Hom}_{\mathrm{Hot}}(Q,R)\\
\pi_0(Q)&\text{---}\ \mathrm{Hom}_{\mathrm{Hot}}(P,R)
\end{aligned}
\]
\deleted{\unsure{donc}}\unsure{en} d'autres termes,
\[
\begin{aligned}
\overset{\deleted{\unsure{\mathrm{Hom}(b}}}{\pi_0(b,c)}&\longrightarrow\mathrm{Hom}_{\mathrm{Hot}}\bigl(\mathrm{Ch}_{\infty}(a,b),\mathrm{Ch}_{\infty}(\unsure{a|b},c)\bigr)\\
\pi_0(a,b)&\longrightarrow\mathrm{Hom}_{\mathrm{Hot}}\bigl(\mathrm{Ch}_{\infty}(b,c),\mathrm{Ch}_{\infty}(a,c)\bigr)
\end{aligned}
\]
\deleted{\unsure{On}} \deleted{\unsure{On en a}} \deleted{\unsure{notera}, pour $\alpha\in\pi_0(b,c)$, par $\alpha*$ \unsure{l'hom.} \unsure{l'homom.} associé de $\mathrm{Ch}_{\infty}(a,b)$ dans $\mathrm{Ch}_{\infty}($} \deleted{Soit $M_{\pi}$ la \unsure{catégorie}} On en déduit en particulier pour $a,b\in M_0$
\[
\left\{\begin{aligned}
1/\ \pi_0(a,a)&\longrightarrow\mathrm{Hom}_{\mathrm{Hot}}\bigl(\mathrm{Ch}_{\infty}(a,b),\mathrm{Ch}_{\infty}(a,b)\bigr)\\
2/\ \pi_0(a,a)&\longrightarrow\mathrm{Hom}_{\deleted{\unsure{\mathrm{Hot}}}\mathrm{Hot}}\bigl(\mathrm{Ch}_{\infty}(b,a),\mathrm{Ch}_{\infty}(b,a)\bigr)
\end{aligned}\right.\tag*{$*)$}
\]
et un objet \inserted{$\alpha_a$ de $\pi_0(a,a)$} est unité bilatère si ses deux images \unsure{soint|soit|sont} les \deleted{foncteurs} \inserted{\unsure{hom.|homs.}} identiques \unsure{de|des} $\mathrm{Ch}_{\infty}(a,b)$ \unsure{resp\textsuperscript{t}} $\mathrm{Ch}_{\infty}(b,a)$. \unsure{Revenant|Passant} aux $\pi_0$ \deleted{\unsure{de}} de ceux-ci, \unsure{cela|cel.} implique que $\alpha$ est unité bilatère pour la composition \deleted{dans $M_{\pi}$ (\unsure{i.e.} \unsure{p.}} des $\pi_0(x,y)$, relativement à $b$. Si on \unsure{veut} que ce soit le cas pour \uline{\unsure{tous} $b$}, \unsure{cela|cel.} signifie \inserted{(prenant $b=a$)} que $\alpha_a$ est un objet unité bilatère pour la structure de comp. associative \inserted{(\unsure{donc} de monoïde)} \deleted{\unsure{dans}} $\pi_0(a,a)$, donc que $\pi_0(a,a)$ est un monoïde \uline{unitaire}, ayant $\alpha$ pour objet unité. L'élément $\alpha_a$ est donc déterminé de façon unique. \deleted{Mais} Et on voit que $\alpha_a$ existe
% archive sheet number: 3

\scanpage{5}\mspage{3}
pour \unsure{tout|tous} $a\in M_0$, donc $M_0$, muni des $\pi_0(a,b)$ et de leur composition, devient une \uline{catégorie}, \deleted{\unsure{de}} soit $M_{\pi}$ ; et les \deleted{\unsure{éléments unité de}} \inserted{flèches identiques} de celle-ci doivent jouer ce rôle d'unité bilatère également pour les ``Hom externe'' $\mathrm{Ch}_{\infty}(a,b)$, à valeurs dans $\mathrm{Hot}\unsure{. |{}_{0}}$

On peut dire que pour $a=a_0$ fixé, les \deleted{$\mathrm{Hom}$} $\deleted{\unsure{\mathrm H}}\mathrm{Ch}_{\infty}(a_0,b)$ ($b$ variable) forment les valeurs d'un \uline{foncteur} covariant
\[
M_{\pi}\unsure{\longrightarrow|\longleftrightarrow}\mathrm{Hot}\qquad b\mapsto\mathrm{Ch}_{\infty}(a_0,b)\unsure{,|;}
\]
\deleted{\unsure{Le fait}} en associant à
\[
b\xrightarrow{u}b'\qquad u\in\pi_0(b,b')
\]
la flèche
\[
\mathrm{Ch}_{\infty}(a,b)\unsure{\text{---}|\longrightarrow}\mathrm{Ch}_{\infty}(b')
\]
\unsure{correspondante} $\unsure{*_1|\alpha_1}$, \unsure{Que|que} ce soit bien là une loi fonctorielle doit résulter formellement de l'associativité des accouplements sur les $\mathrm{Ch}_{\infty}$. De $\unsure{\overline{\mathrm m}|=}$, on a une loi fonctorielle, pour $b$ fixé
\[
M_{\pi}^{o}\longrightarrow\mathrm{Hot}\qquad a\unsure{\vdash|\mapsto}\mathrm{Ch}_{\infty}(a,b),
\]
grâce à l'associativité. En fait, l'associativité impliquera aussi qu'on a un \uline{bifoncteur}
\[
\begin{aligned}
M_{\pi}^{o}\times M_{\pi}&\longrightarrow\mathrm{Hot}\\
(a,b)&\longrightarrow\deleted{\unsure{\mathrm H}}\mathrm{Ch}_{\infty}(a,b),
\end{aligned}
\]
% archive sheet number: 4

\scanpage{6}\mspage{4}
car si
\[
a'\overset{u}{\unsure{\text{---}|\longrightarrow}}a\qquad b\xrightarrow{v}b'
\]
on a le diagramme
\diagram{\begin{tikzcd}[ampersand replacement=\&,column sep=large,row sep=large]
\mathrm{Ch}_{\infty}(a,b) \arrow[r,no head,"{\substack{\text{comp. avec}\\ v}}"] \arrow[d,"{\text{comp. avec}}"',"{u}"] \& \mathrm{Ch}_{\infty}(a,b') \arrow[d,"{\text{comp. avec }u}"]\\
\mathrm{Ch}_{\infty}(a',b) \arrow[r,no head,"{\substack{\text{comp. avec}\\v}}"] \& \mathrm{Ch}_{\infty}(a',b')
\end{tikzcd}}
qui est commutatif grâce à l'associativité pour
\[
\mathrm{Ch}_{\infty}(a',a),\ \mathrm{Ch}_{\infty}(a,b),\ \mathrm{Ch}(b,b').
\]

\uline{Corollaire} de cette \unsure{dépendance|disposition} bifonctorielle : Si $u:\deleted{\unsure{b}}a'\to\deleted{\unsure{b'}}a$, \inserted{$v:b\to b'$} \unsure{st|est} \unsure{isom.} dans $M_{\pi}$, alors la flèche correspondante
\[
\mathrm{Ch}_{\infty}(a,b)\longrightarrow\mathrm{Ch}_{\infty}(a',b')\qquad\text{dans Hot}
\]
est un \uline{iso}.

Les $\unsure{\overline{\mathrm m}\mathrm s|\overline{\mathrm m}\mathrm e\mathrm s}$ arguments devraient marcher si au lieu \unsure{des|du} $\mathrm{Hom}$ ext. à valeurs dans $\mathrm{Hot}$, on en prenait dans $\mathrm{Ind}\,\mathrm{Hot}$, ou dans $\mathrm{Cat}\,\mathrm{hot}$, ou dans $\mathrm{Ind}\,\mathrm{Cat}\,\mathrm{hot}$ -- je pense bien sûr aux $\mathrm{CH}_{\infty}(a,b)$ (dans la théorie de la localisation) dans $\mathrm{Ind}\,\mathrm{Cat}\,\mathrm{hot}$, qui \deleted{définissent} peuvent être regardés comme provenant d'un bifoncteur (où $M_{\pi}=M\Sigma^{-1}$)
\[
M_{\pi}^{\unsure{o}}\times M_{\pi}\ \text{---}\ \mathrm{Ind}\,\mathrm{Cat}\,\mathrm{hot}
\]
% archive sheet number: 5

\scanpage{7}\mspage{5}
Voici l'axiomatisation définitive (?). On remplace $\mathrm{Hot}$ par une catégorie
% slanted
\margin{On ne doit pas \unsure{avoir} besoin de la commutativité}
$\unsure{\mathscr H|\mathcal H}$, catégorie avec produit $*$ associatif, \deleted{\unsure{et}} commutatif (?), unitaire, et toutes les compatibilités habituelles \unsure{au} besoin,
l'objet unité pour $*$ sera noté $1$, et on notera
\[
\left\{\begin{aligned}
\pi_0(\unsure{\xi|\zeta})&\overset{\text{déf}}{=}\mathrm{Hom}_{\mathscr H}(1,\xi)\\
\pi_0:\mathscr H&\longrightarrow\mathrm{Ens}
\end{aligned}\right.
\]

Si on a
\[
P*Q\xrightarrow{u}R
\]
on en déduit
\[
\left\{\begin{aligned}
\pi_0(P)&\overset{c_u^{\unsure{\ell|l}}}{\longrightarrow}\mathrm{Hom}_{\mathscr H}(Q,R)\\
\pi_0(Q)&\overset{c_u^r}{\longrightarrow}\mathrm{Hom}_{\mathscr H}(P,R)
\end{aligned}\right.
\]
ainsi :
% Upright sign in the left margin of the following formula.
\margin{\unsure{!|]}}
\[
\begin{aligned}
&\text{Si }\alpha\in\pi_0(P)\ \text{i.e.}\ 1\xrightarrow{\alpha}P,\ \text{alors }\alpha*\mathrm{id}_Q:Q\to P*Q\\
&\hspace{4em}\text{et }\deleted{\unsure{c(\alpha)}}\ c_u^\ell(\alpha)=u\circ(\alpha*\mathrm{id}_Q)\\
&\text{et dualement}\qquad c_u^r(\beta)=u\circ(\mathrm{id}_{\unsure{p|P}}*\beta)
\end{aligned}
\]
\deleted{\unsure{On}} Si on a trois \deleted{\unsure{ensembles}} \inserted{\unsure{objets}} $P_0,P_1,P_2$ dans $\mathscr H$ des accouplements
\diagram{\begin{tikzcd}[ampersand replacement=\&,column sep=small]
a \arrow[r,"{P_0}"',no head] \arrow[rr,bend left=23,"{Q_0}",no head] \arrow[rrr,bend left=38,"R",no head] \& b \arrow[r,"{P_1}"',no head] \arrow[rr,bend right=28,"{Q_1}"',no head] \& c \arrow[r,no head,"{P_2}"'] \& d
\end{tikzcd}}
\[
\left\{\begin{aligned}
P_0*P_1&\xrightarrow{u_1}Q_0 & P_1*\deleted{\unsure{Q}}P_2&\overset{u_2}{\unsure{\text{---}|\longrightarrow}}Q_1\\
&\deleted{Q_0*Q_1}\\
P_0*Q_1&\xrightarrow{u_3}R\\
Q_0*P_2&\overset{u_4}{\text{---}}R
\end{aligned}\right.
\]
% archive sheet number: 6

\scanpage{8}\mspage{6}
avec une relation d'associativité, disant que
\diagram{\begin{tikzcd}[ampersand replacement=\&]
P_0*P_{\unsure{1|3}}*P_2 \arrow[r] \arrow[d] \& Q_0*P_2 \arrow[d]\\
P_0*Q_1 \arrow[r] \& R
\end{tikzcd}}
est commutatif, alors on en déduit \deleted{trois} \inserted{deux} relations de transitivité et une de commutation, pour les opérations associées aux éléments de $\pi_0(P_0),\pi_0(P_1),\pi_0(P_2)$, qu'on \unsure{n'explicitera} pas, sauf l'une :
\[
\text{Soit }\beta\in\deleted{\unsure{\pi_0}}P_1,\ \gamma\unsure{\in|:}P_2
\]
\[
\begin{aligned}
\text{d'où par }\beta:\quad &c_{u_1}(\beta):P_0\longrightarrow Q_0\\
\text{et par }\gamma:\quad &c_{u_4}(\gamma):Q_0\longrightarrow R\\
\text{et par }\underset{\substack{\rotatebox{-90}{$\in$}\\\pi_0(Q_1)}}{\deleted{\unsure{\beta\gamma}}\inserted{\beta\underset{u_2}{*}\gamma}}:\quad &c_{u_3}(\beta\gamma):P_0\longrightarrow\deleted{\unsure{Q}}R
\end{aligned}
\]
\unsure{eh|et} bien on a
\[
\boxed{c(\beta\underset{u_2}{*}\gamma)=c(\gamma)c(\beta)}.
\]
Il faut encore définir $\beta\underset{u_2}{*}\gamma$, via l'accouplement $P_1*P_2\xrightarrow{u_2}Q_1$, on a en effet
\[
\begin{aligned}
\pi_0(P_1)\times\pi_0(P_{\unsure{7|2|1}})&\longrightarrow\pi_0(P_1*P_2)\\
(\beta,\gamma)&\longmapsto\beta*\gamma
\end{aligned}
\]
et on compose avec
\[
\pi_0(u_2):P_1*P_2\longrightarrow Q_1\qquad\unsure{?|:}
\]
% archive sheet number: 7

\scanpage{9}\mspage{7}
On se donne un ens. $M_0$, et
\[
\left\{\begin{aligned}
M_0\times M_0&\overset{\mathrm{Ch}}{\longrightarrow}\unsure{\mathscr H|\mathcal H}&\text{noté }\quad(a,b)&\longrightarrow\mathrm{Ch}(a,b)\\
&&&\text{ou }\mathrm{Hom}\,\mathrm{ext}_{\mathscr H}(a,b)\\
M_0\times M_0\times M_0&\overset{\mathrm{comp}}{\longrightarrow}\mathrm{Fl}\,\mathscr H &(a,b,c)&\mapsto\bigl(\mathrm{Ch}(a,b)\times\mathrm{Ch}(b,c)\\
&&&\hspace{3em}\longrightarrow\mathrm{Ch}(a,c)\bigr)
\end{aligned}\right.
\]
avec deux axiomes

a) Associativité pour les compositions, pour $(a,b,c,d)\in M_0^{\unsure{4}}$
\[
\begin{aligned}
&\deleted{\unsure{M_0\times M_0\times M_0\times M_0}\longrightarrow}\\
&\deleted{\mathrm{Ch}(a,b)\times\mathrm{Ch}(b,c)\times\mathrm{Ch}(c,d)\longrightarrow\mathrm{Ch}(a,d)}
\end{aligned}
\]
\diagram{\begin{tikzcd}[ampersand replacement=\&,column sep=small]
\mathrm{Ch}(a,b)\times\mathrm{Ch}(b,c)\times\mathrm{Ch}(c,d) \arrow[r] \arrow[d] \& \mathrm{Ch}(a,c)\times\mathrm{Ch}(c,d) \arrow[d]\\
\mathrm{Ch}(a,b)\times\mathrm{Ch}(b,d) \arrow[r] \& \mathrm{Ch}(a,d)
\end{tikzcd}}

b) Existence d'objets unité bilatères
\[
1_a\in\pi_0(\mathrm{Ch}(a,a))\qquad\text{pour tt }a\in M_0.
\]
Alors les $\pi_0(\mathrm{Ch}(a,b))$ et leur composition définissent sur $M_0$\deleted{\unsure{,}} \inserted{comme ens. d'objets} une structure de \uline{catégorie} \inserted{$M_{0\pi}$}, ayant les $1_a$ comme identités, et on trouve que
\[
(a,b)\unsure{\vdash|\mapsto}\mathrm{Ch}(a,b)
\]
peut s'interpréter \unsure{à partir|: partir} d'un bi\uline{foncteur} canonique
\[
M_{0\pi}^{o}\times M_{0\pi}\longrightarrow\mathscr H
\]
dont le composé avec
\[
\mathscr H\xrightarrow{\pi_0}\mathrm{Ens}
\]
n'est autre que le bifoncteur $(a,b)\longrightarrow\mathrm{Hom}_{M_{0\pi}}(a,b)$,
% archive sheet number: 8

\scanpage{10}\mspage{8}
\[
\mathrm{Hom}_{M_{0\pi}}(a,b)\simeq\pi_0\bigl(\mathrm{Ch}(a,b)\bigr).
\]
\deleted{\unsure{Cf. cette}} Cette relation est la \uline{définition} de $\mathrm{Hom}_{M_{0\pi}}$, mais ici la relation est entendue comme une bijection \uline{fonctorielle} en $a,b$.

On trouve notamment que si
\[
u:a'\unsure{\text{---}|\longrightarrow}a,\qquad v:b\ \text{---}\ b'
\]
sont des flèches \uline{inversibles} \inserted{\unsure{de|dans} $M_{0\pi}$}, alors
\[
\mathrm{Ch}_{\deleted{\unsure{\infty}}}(a,b)\longrightarrow\mathrm{Ch}(a',b')
\]
correspondant dans $\unsure{\mathscr H|\mathcal H}$ est \uline{iso}.

On appliquera tout ceci \unsure{surtout} \unsure{aux} deux situations déduites d'une situation $(M,\Sigma)$, avec
\[
\begin{aligned}
1^{\mathrm o})\quad &\mathscr H=\mathrm{Hot},&
 \mathrm{Ch}(a,b)&=\mathrm{Ch}^{\Sigma}_{\infty\deleted{\unsure{\Sigma}}}(M;a,b)\\
2^{\mathrm o})\quad &\mathscr H=\inserted{\mathrm{Ind}}\,\mathrm{Cat}\,\mathrm{hot},&
 \mathrm{Ch}(a,b)&=\mathrm{CH}^{\Sigma}_{\infty}(M;a,b).
\end{aligned}
\]
% slanted
\margin{$*$ dans $\mathrm{Hot}$ et $\mathrm{Ind}\,\mathrm{Cat}\,\mathrm{hot}$ est le produit cartésien sans plus}
Dans les deux cas, on a donc \deleted{\unsure{un}} un foncteur sur $(M\Sigma^{-1})^{o}\times(M\Sigma^{-1})$, à valeurs soit dans $\mathrm{Hot}$, soit dans \deleted{$\unsure{\mathrm{Cat}}$} $\mathrm{Ind}\,\mathrm{Cat}\,\mathrm{hot}$, (\uline{NB} les deux structures de cat. sur $M_0$ déduites de ces deux foncteurs sont can\textsuperscript{t} isomorphes)
\diagram{\begin{tikzcd}[ampersand replacement=\&,row sep=small,column sep=small]
\& \mathrm{Hot} \& (a,b)\unsure{\longrightarrow|\twoheadrightarrow}\mathrm{Ch}_{\infty}^{\Sigma}(a,b)\\
(M\Sigma^{-1})^o\times(M\Sigma^{-1}) \arrow[ur] \arrow[dr] \& \& \\
\& \mathrm{Ind}\,\mathrm{Cat}\,\mathrm{hot} \& (a,b)\ \text{---}\ \mathrm{CH}_{\infty}^{\Sigma}(a,b)\unsure{,|.}
\end{tikzcd}}
% archive sheet number: 9

\scanpage{11}\mspage{9}
Comme les foncteurs canoniques
\diagram{\begin{tikzcd}[ampersand replacement=\&,column sep=small]
\mathrm{Ind}\,\mathrm{Cat}\,\mathrm{hot} \arrow[r,"{\mathrm{Ind}\,\unsure{\varphi|\psi|\phi}}"] \& \mathrm{Ind}\,\mathrm{Hot} \arrow[r,"{\mathrm{Ind}\,\unsure{\tau|+|\gamma}}"] \& \mathrm{Ind}\,\mathrm{Hot}_0^{\wedge} \arrow[r,"{\varinjlim}"] \& \mathrm{Hot}_0^{\wedge}\\
\mathrm{Cat}\,\mathrm{hot} \arrow[u,hook] \arrow[r,"{\varphi}"] \& \mathrm{Hot} \arrow[u,hook] \arrow[r,"{\tau}"] \& \mathrm{Hot}_0^{\wedge} \arrow[u,hook] \arrow[ur,"{\mathrm{id}}"'] \&
\end{tikzcd}}
commutent aux produits finis, toute théorie \unsure{de|des} ``Hom externes'' à valeurs dans une de ces catégories, en définit une à valeurs dans chacune de celles dans lesquelles elle s'envoie.
\deleted{\unsure{Aussi}} D'ailleurs, ces foncteurs \deleted{sont} ne sont pas tous pl. fid. (c'est en tous cas faux pour $\varphi$ et $\mathrm{Ind}\,\varphi$, et j'ignore ce qu'il en est de $\tau$ et $\mathrm{Ind}\,\tau$ -- les flèches verticales par contre sont pl. fid.)
% slanted
\margin{et \unsure{églmt|égt} $\mathrm{Ind}\,\mathrm{Hot}_0^{\wedge}\xrightarrow{\varinjlim}\mathrm{Hot}_0^{\wedge}$ n'est \uline{pas} pl. fidèle}
Mais ils induisent des bijections pour les $\pi_0$. Donc la catégorie $M_{0\pi}$ déduite d'une théorie des Hom externes \inserted{à val.} dans une de ces catégories, et celle déduite dans une autre où elle s'envoie, est la $=$ à isom. can. près. 
Pour ce qui est de cette définition \inserted{de $M_{0\pi}$}, rien n'est changé donc quand on regarde les $\mathrm{CH}_{\infty}$ comme des objets de $\mathrm{Ind}\,\mathrm{Hot}$, ou seulement de $\mathrm{Ind}\,\mathrm{Hot}_0^{\wedge}$\inserted{ou $=$ \deleted{\unsure{slt}} de $\mathrm{Hot}_0^{\wedge}$}, ou quand on regarde les $\mathrm{Ch}_{\infty}$ comme des objets
% archive sheet number: 10

\scanpage{12}\mspage{10}
\unsure{de|dans} \deleted{$\unsure{\mathrm{Ch}_{\infty}}$} $\mathrm{Hot}_0^{\wedge}$. Mais \deleted{\unsure{de}}\unsure{bien} sûr, la donnée des accouplements de composition perd chaque fois en précision (sauf peut-être en passant de $\mathrm{Hot}$ à $\mathrm{Hot}_0^{\wedge}$, si $\unsure{\tau|+|\gamma}$ est pl. fidèle \ldots).

Quand on prend les images des $\mathrm{CH}_{\infty}^{\Sigma}(a,b)$ et $\mathrm{Ch}_{\infty}^{\Sigma}(a,b)$ dans $\mathrm{Hot}_0^{\wedge}$ (la dernière catégorie du diagramme), on trouve des objets can. isomorphes (fonctoriellement en $a,b$ dans $M_{0\pi}\simeq M\Sigma^{-1}$). Quand on les regarde comme étant resp. dans $\mathrm{Ind}\,\mathrm{Hot}$ et dans $\mathrm{Hot}$ \inserted{$\hookrightarrow\mathrm{Ind}\,\mathrm{Hot}$}, on trouve un \unsure{hom|hom.} fonctoriel
\[
\mathrm{CH}_{\infty}^{\Sigma}(a,b)\longrightarrow\mathrm{Ch}_{\infty}^{\Sigma}(a,b)
\qquad\begin{aligned}&\text{dans}\\&\mathrm{Ind}\,\mathrm{Hot},\\&\text{ou dans}\\&\mathrm{Ind}\,\mathrm{Hot}_0^{\wedge}\\&\text{au choix}\end{aligned}
\tag*{($\alpha$)}
\]
\unsure{Les|La} $\varinjlim$ de $\mathrm{CH}_{\infty}^{\Sigma}(a,b)$ dans $\mathrm{Hot}$ n'existe pas en général (sauf si le $\mathrm{Ind}$-objet est \unsure{ess\textsuperscript{t}} constant), mais \unsure{il|elle} existe dans \deleted{$\mathrm{Ind}$} $\mathrm{Hot}_0^{\wedge}$, et l'isom. \unsure{can|can.}
\[
\underset{\mathrm{Hot}_0^{\wedge}}{\varinjlim}\mathrm{CH}_{\infty}^{\Sigma}(a,b)\simeq\mathrm{Ch}_{\infty}^{\Sigma}(a\unsure{\mid|,|;}b)\quad\text{dans }\mathrm{Hot}_0^{\wedge}
\tag*{($\alpha\alpha$)}
\]
est la \deleted{\unsure{classe}} \inserted{flèche} déduite de $(\alpha)$ dans $\mathrm{Ind}\,\mathrm{Hot}_0^{\wedge}$,
% Isolated ink at the foot, visibly separate from the last line.
\figure{\begin{tikzpicture}[x=1mm,y=1mm,line cap=round,line join=round]
\draw[line width=.16mm] (.65,2.52) .. controls (1.36,2.44) and (2.13,1.91) .. (2.66,2.06);
\fill (.32,1.95) ellipse (.23mm and .15mm);
\draw[line width=.18mm] (.97,1.82) .. controls (1.48,1.68) and (2.03,1.34) .. (2.36,1.43);
\fill (3.02,2.34) ellipse (.53mm and .49mm);
\draw[line width=.35mm] (3.14,2.18) .. controls (3.13,1.48) and (2.96,.93) .. (2.76,.61) .. controls (2.2,.33) and (1.03,.62) .. (.54,.49);
\draw[line width=.23mm] (3.09,.32) .. controls (3.61,.18) and (4.04,.22) .. (4.36,.32);
\fill (4.43,.31) ellipse (.24mm and .16mm);
\draw[line width=.16mm] (.88,.04) .. controls (1.31,-.04) and (1.77,-.26) .. (2.36,-.32);
\end{tikzpicture}}
% archive sheet number: 11

\scanpage{13}\mspage{11}
Dans les deux cas $\mathrm{CH}$ et $\mathrm{Ch}$ (qu'on va désigner par le symbole commun $C_{\infty}(a,b)$), la fonctorialité de cet objet $C_{\infty}(a,b)$ (\unsure{$\overline{\mathrm m}$ême|même} \unsure{au|en} sens le plus faible, comme objet de $\mathrm{Hot}_0^{\wedge}$) par rapport à $a$ et $b$ ($=$ au sens le plus faible, comme foncteur $M^o\times M\longrightarrow\mathrm{Hot}_0^{\wedge}$ et pas même $(M\Sigma^{-1})^o\times M\Sigma^{-1}\to\unsure{\mathrm{Hot}_0^{\wedge}|\mathrm{Hot}}$) n'est pas évidente du tout sur la définition directe \deleted{\unsure{des}} en termes du système inductif explicite dans $\mathrm{Cat}$
\[
\unsure{\mathbb C|C}_{\infty}(a,b)=\bigl(\mathrm{Ch}_{\tau_0^{2n}}(a,b)\bigr)_{n\in\mathbb N},
\]
car une flèche $a'\longrightarrow\unsure{a}$ ou $b\unsure{\text{---}|\longrightarrow}b'$ dans $M$, et même qu'elle serait dans $\Sigma$, ne définit aucunement une \deleted{\unsure{application}} flèche dans $\mathrm{Cat}$
\[
\begin{aligned}
C^{\deleted{\unsure{h|\Sigma}}}_{\infty}(a,b)&\longrightarrow C_{\infty}(a',b)\\
\text{ou}\qquad C_{\infty}(a,b)&\ \text{---}\ C_{\infty}(a,b').
\end{aligned}
\]
Cette difficulté rappelle celle qu'avait à surmonter Quillen, dans Homotopical Algebra, pour affronter la différence entre ``left homotopies'' et ``right homotopies''. Je me rappelle comment
% archive sheet number: 12

\scanpage{14}\mspage{12}
on peut néanmoins définir de telles flèches, comme résultant d'une \deleted{application} \inserted{flèche} (beaucoup plus riche) \deleted{de} ``composition'' de chemins entre $a'$ et $a$, et $a$ et $b$, ou entre $a$ et $b$, et $b$ et $b'$, soit en général donc l'accouplement
\[
C_{\infty}(a,b)\times C_{\infty}(b,c)\longrightarrow C_{\infty}(a,c),
\]
dont la définition directe n'est pas moins problématique, et \unsure{on|ne} peut en tous cas \uline{pas} le \deleted{faire} \inserted{donner} comme une flèche dans $\mathrm{Cat}$ \unsure{tout-même|elle-même|tout-mieux}, ou dans $\mathrm{Ind}\,\mathrm{Cat}$.

\deleted{\unsure{Pour y}} Interprétant \unsure{p.ex} \deleted{\unsure{des chemins}} \inserted{\unsure{les|ces}} $\mathrm{Ch}_{\infty}(a,b)$ en termes de chemins \uline{infinis} (des deux côtés) de $a$ à $b$, en un sens évident (mais ``essentiellement finis'' des deux côtés), on peut définir aussi
\[
\mathrm{Ch}_{\infty}^{\deleted{\unsure{+|\Sigma}}-}(a,b),\quad\mathrm{Ch}_{\infty}^{+}(a,b)
\]
comme les sous-catégories pleines formées \deleted{des} respectivement de chemins qui sont \inserted{\unsure{t.t.|tôt}} dégénérés en degrés $\geq0$ (i.e. $a_0\leftarrow a_1\to a_2\leftarrow a_3\ldots$ sont tous des identités), ce qui revient aux chemins infinis seulement à gauche, i.e. s'arrêtant à $a_0$, et symétriquement pour $\mathrm{Ch}_{\infty}^{+}(a,b)$. Même \unsure{tabac|chose|topo} pour les $\mathrm{CH}^{+},\mathrm{CH}^{-}_{\unsure{\infty}}$.
% archive sheet number: 13

\scanpage{15}\mspage{13}
Ceci dit, on a ceci (revenant à la notation $C_{\infty}$ ambivalente) : \inserted{a)} les inclusions \unsure{ple|pl.} fid.
\diagram{\begin{tikzcd}[ampersand replacement=\&]
\unsure{\text{a)}}\quad C_{\infty}^{\deleted{\unsure{-|+}}}(a,b) \arrow[r,hook] \& C_{\infty}^{\deleted{\unsure{-|+}}}(a,b) \& C_{\infty}^{+}(a,b) \arrow[l,hook']
\end{tikzcd}}
sont des iso dans $\mathrm{Ind}\,\mathrm{Cat}\,\mathrm{hot}$ resp. dans $\mathrm{Hot}$, et

b) On a bel et bien \unsure{une} composition des chemins
\[
C_{\infty}^{-}(a,b)\times C_{\infty}^{+}(b,c)\longrightarrow C_{\infty}(a,c).
\]
L'\uline{associativité} dans le cas \unsure{du|des} $\mathrm{CH}$ n'offre d'ailleurs pas de difficulté, à cause de la possibilité de se \deleted{\unsure{laisser}} \unsure{ramener|revenir} aux
\[
C\deleted{\unsure{H}}_{\unsure{m|m_1}}(a,b)\times C\deleted{\unsure{H}}_n(b,c)\times C_p(c,d)\rightrightarrows C_N(a,d)
\]
pour $N$ grand, et prouver l'égalité \deleted{\unsure{dans C}} comme flèches dans $\mathrm{Cat}\ \mathrm{hot}$ -- ce qui résultera des lemmes d'homotopie sur les intervalles dans $\mathrm{Cat}$. Dans le cas des $\mathrm{Ch}_{\infty}$ par contre je n'ai pas trouvé de démonstration. (Faute seulement d'avoir essayé sérieusement de m'y mettre ? Je n'ai pas trop envie, vu qu'il se pourrait que $\mathrm{Hot}\unsure{\text{---}|\to}\mathrm{Hot}_0^{\wedge}$ soit pleinement fidèle, auquel cas le travail devient \unsure{inutile} \ldots)
% archive sheet number: 14

\scanpage{16}\mspage{14}
2. \uline{Comparaison avec} \uline{Quillen}.

\deleted{\unsure{On}} Il se pose un certain nombre de questions sur ce foncteur $(M\Sigma^{-1})^o\times M\Sigma^{-1}\longrightarrow\mathrm{Hot}$, et sa relation aux constructions classiques, notamment à la théorie de Quillen. Je vais maintenant noter
\[
\begin{aligned}
\mathrm{Ch}_{\infty}^{\Sigma}(a,b)&=\mathrm{Hom}\,\mathrm{ext}_{\mathrm{Hot}}(a,b)\quad\text{ou simplement}\\
&\quad\mathrm{Hom}\,\mathrm{ext}(a,b).
\end{aligned}
\]

$1^{\mathrm o})$ Soient $f,g:a\rightrightarrows b$ dans $M$, alors (si $M$ est une ``model category'') Quillen définit un ensemble
\[
\pi_1(a,b;f,g)\quad\text{ou}\quad\pi_1(f\unsure{,|.}g)
\]
des classes d'homotopie de $f$ à \unsure{$g$}, de telle façon que $\mathrm{Hom}_{M}(a,b)$ devient \inserted{l'\unsure{ens.} des obj. d'} un groupoïde, dont les flèches sont les éléments des $\pi_1(f,g)$, pour une notion de composition naturelle.

Je suspecte que l'on a
\[
\pi_1(f,g)=\mathrm{Hom}_{\underbrace{\Pi(\mathrm{C}\underline{\mathrm{h}}_{\infty}^{\Sigma}(a,b))}_{\substack{\text{groupoïde}\\\text{fondamental}\\\text{de }\mathrm{C}\underline{\mathrm{h}}_{\infty}^{\Sigma}(a,b)}}}(f,g),
\]
quand $f$ et $g$ sont considérés comme des ``chemins infinis'' en \unsure{inscrivant|insérant} chaque $f:a\deleted{\unsure{\leftarrow}}\unsure{\text{---}|\to}b$ en
\diagram{\begin{tikzcd}[ampersand replacement=\&,column sep=small]
\ldots \& \underset{\substack{\Vert\\a}}{a_{-2}} \& \underset{\substack{\Vert\\a}}{a_{-1}} \arrow[l,"{\mathrm{id}}"'] \arrow[r,"f"] \& \underset{\substack{\Vert\\b}}{a_0} \& \underset{\substack{\Vert\\b}}{a_1} \arrow[l,"{\mathrm{id}}"'] \arrow[r] \& \ldots
\end{tikzcd}}
% archive sheet number: 15

\scanpage{17}\mspage{15}
Donc la définition des $\pi_1(f,g)$ \inserted{(si c'est bien \unsure{ça|sa})} se prolonge de façon naturelle aux chemins infinis (\unsure{ess\textsuperscript{t}|ess.} finis) arbitraires.

$2^{\mathrm o})$ Supposons que $M$ admette un objet à la fois initial et final $0_M$ ($M$ ``\unsure{ponctuée|pointée}''). Je dis qu'alors \deleted{\unsure{toutes}} les catégories $\underline{\mathrm{Ch}}_{\infty}^{\Sigma}(a,b)$ sont \unsure{ponctuées|pointées} de façon naturelle : à savoir, par le chemin infini de $a$ à $b$ défini \inserted{(comme plus haut)} à partir de $f=0_{a,b}:a\unsure{\text{---}|\longrightarrow}b$, l'application nulle de $a$ à $b$, On \deleted{peut} a envie alors de regarder les $\underline{\mathrm{Ch}}_{\infty}^{\Sigma}(a,b)$ comme définissant des éléments dans le $\mathrm{Hot}_{\bullet}$ (``Hot \unsure{ponctué|pointé}''). C'est plus fin que de dire qu'on se donne une classe neutre
\[
e\longrightarrow\mathrm{Ch}(a,b)\qquad\text{dans Hot}
\]
car le foncteur de $\mathrm{Hot}_{\bullet}$ dans $\mathrm{Hot}$ (ou dans $e\backslash\mathrm{Hot}$) n'est pas \unsure{fidèle}\unsure{?|.} Le problème, c'est de définir les flèches
\[
\mathrm{Ch}(a,b)\times\mathrm{Ch}(b,c)\ \text{---}\ \mathrm{Ch}(a,c)
\]
dans $\mathrm{Hot}_{\bullet}$, et pas seulement dans $\mathrm{Hot}$, et avec associativité \inserted{(et les unités)} dans $\mathrm{Hot}_{\bullet}$ également. Si on applique la définition avec $\mathrm{Ch}^{-}(\unsure{a},b)$ et $\mathrm{Ch}^{+}(b,c)$, le composé des éléments nuls
\[
\begin{aligned}
&\unsure{a}\overset{0}{\text{---}}b\ \deleted{\unsure{\longleftarrow b\longrightarrow c}}\quad\text{de }\mathrm{Ch}^{\deleted{\unsure{-|+}}}(a,b)\quad\text{et}\\
&\deleted{\unsure{b\longrightarrow b\longrightarrow c}}\inserted{b\xleftarrow{\mathrm{id}}b\xrightarrow{0}c}\quad\text{de }\mathrm{Ch}^{+}(b,c)
\end{aligned}
\]
% archive sheet number: 16

